% ============================================================================
% SECTION 5 --- Experiments.  Order: Correctness -> Competence -> Attribution
%               -> Capabilities.  Declarative subsection titles.
%
% EVERY NUMBER HERE CARRIES A "% source:" COMMENT.  The verified set is the
% seven files listed at the top of main.tex and nothing else.
% EVERY UNMEASURED QUANTITY IS \XXX INSIDE A "GAP" BLOCK.  Do not replace an
% \XXX with a plausible value; replace it with a measurement or leave it.
% Compressed to the ICLR 9-page limit.  Full parity list -> app:audit; full
% experimental designs -> app:capdesign; the 2-objective table ->
% app:multiobjtable.  The GAP blocks are load-bearing and must NOT be
% shortened to save space.
% ============================================================================
\section{Experiments}
\label{sec:experiments}

We evaluate \method{} from the learned molecular reference process through finite-horizon control, standard molecular editing, multi-objective design, and trajectory-level intervention.

\subsection{\method{} learns state-dependent preferences among executable molecular edits}
\label{sec:exp-rtheta}

The executable rewrite system determines which molecular edits are possible, but not which should be preferred in a given molecular state. We therefore test whether the frozen reference law $R_\theta$ learns state-dependent transition structure beyond corpus-wide edit-family frequencies. On $3{,}545$ held-out transitions with source molecules disjoint from reference-law training, we compare $R_\theta$ with an empirical-family baseline that uses training-set family frequencies while operating over the same canonical molecular successors. This holds executable support fixed and asks only whether learning improves how probability is distributed among the available next molecules.

$R_\theta$ reduces canonical-successor NLL from $5.37$ to $3.90$ nats relative to the empirical-family baseline (\cref{fig:mechanism}(a)). A matched decomposition shows that $86.7\%$ of this improvement is recovered by retaining empirical family frequencies and learning only which successor is preferred within the current state. Thus most of the learned signal is not which kinds of edits occur globally, but which executable molecular edit is appropriate for the molecule at hand. The rewrite system therefore defines what is possible, while $R_\theta$ supplies a learned, state-dependent preference over those possibilities.

This effect persists when evaluation is restricted to edit families with data-backed transition supervision, where $R_\theta$ retains a $1.06$-nat advantage over empirical family frequencies (\cref{app:exp-rtheta}). The stronger gains observed for synthetically supervised ring-edit families are therefore not required for the state-dependent preference claim. We use $R_\theta$ throughout the remaining experiments as a goal-independent reference law over fixed executable support, without interpreting its probabilities as physical kinetics or synthetic-route likelihood.

\subsection{\method{} uses future reachability to improve multi-step molecular editing}
\label{sec:exp-lookahead}

\method{} can evaluate a molecular edit by the futures it leaves reachable, rather than only by the molecule it produces immediately. Because control acts on canonical molecular successors and every committed successor is itself a complete molecule, each candidate edit defines a realized molecular intermediate from which the frozen reference process can be continued over the remaining budget. We isolate this capability through exact target recovery on $65$ held-out source--target pairs with source molecules disjoint from development, each admitting a certified executable edit path within budget. Both policies operate over the same frozen $R_\theta$, target, edit budget, and candidate successors. The local policy greedily selects the successor most similar to the target, whereas the future-aware rollout policy evaluates the continuations reachable over the remaining budget and overrides the local choice only when doing so yields a strict improvement. We independently verify the corresponding exact finite-horizon construction on an enumerable molecular state graph (\cref{app:exp-lookahead}).

Accounting for future reachability changes which molecular edits succeed. Remaining-budget rollout recovers $40/65$ targets, compared with $26/65$ under local selection, rescuing $14$ of the $39$ local failures ($+21.5$ percentage points, $[12.3,33.5]$, \cref{fig:mechanism}(b--c)). A deterministically selected rescue in \cref{fig:mechanism}(b) illustrates the mechanism. The local policy chooses the higher-similarity successor, whereas a single future-aware override takes a lower-similarity edit that preserves a successful continuation and reaches the target within budget. Across the full panel, these rescues show that the value of an edit depends not only on the molecule it produces immediately, but on the molecular futures it leaves reachable.

At molecular scale, exhaustive evaluation of these continuations is expensive, so we use $h_\phi$ to prioritize which successors receive rollout evaluation. This retains the same aggregate recovery using only $34.6\%$ of the continuation evaluations required by all-candidate lookahead (\cref{fig:mechanism}(d)). We treat this result as a prioritization diagnostic rather than a standalone evaluation of $h_\phi$, because the two procedures do not recover exactly the same targets, and exact-target similarity uses information unavailable in ordinary property optimization. \Cref{sec:exp-pareto} removes this privileged target information and tests whether the same future-reachability principle improves target-free multi-objective molecular design.

%% TITLE IS TEMPORARY: 'performs strongly' until the official 800xK=20 cohort
%% licenses 'is competitive'.  One word changes; nothing else.
\subsection{\method{} performs strongly on standard source-conditioned molecular editing}
\label{sec:exp-jin}

We next ask whether the complete \method{} system is competitive on a standard molecular-editing task. We use the ZINC-250k similarity-constrained QED benchmark reported by GrIDDD \citep{ninniri2025griddd}. Source molecules begin with QED in $[0.7,0.8]$, and a source is solved if at least one returned molecule reaches $\mathrm{QED}\ge0.9$ while retaining Morgan-fingerprint Tanimoto similarity of at least $0.4$. The reference law $R_\theta$ is unchanged from \cref{sec:exp-rtheta,sec:exp-lookahead}, and only the QED future-value controller is trained for this task, using sources disjoint from evaluation. Each candidate slot corresponds to one controlled trajectory. Successful trajectories terminate prospectively upon first entering the target region, while trajectories that never reach the region remain failures rather than being resampled.

\begin{table}[!ht]
\centering
{\small
\begin{tabular}{lrrrr}
\toprule
Method & Sources & $K$ & QED success ($\uparrow$) & Diversity ($\uparrow$) \\
\midrule
JT-VAE \citep{jin2018jtvae}   & 800 & 20 & 8.8\%  & -- \\
CG-VAE \citep{liu2018cgvae}   & 800 & 20 & 4.8\%  & -- \\
GCPN \citep{you2018gcpn}      & 800 & 20 & 9.4\%  & 0.216 \\
GrIDDD \citep{ninniri2025griddd} & 800 & 20 & 45.1\% & 0.283$^{\dagger}$ \\
\midrule
\textbf{\method{}} & 128$^{*}$ & 12$^{*}$ & 54.7\%$^{*}$ & 0.393$^{*}$ \\
\bottomrule
\end{tabular}}
\caption{\small\textbf{Source-conditioned editing on the Jin/ZINC-250k similarity-constrained QED benchmark.} Published baseline values are reproduced from the comparator set reported by GrIDDD. $K$ denotes candidates returned per source, and success requires $\mathrm{QED}\ge0.9$ and Tanimoto similarity $\ge0.4$. $^{*}$\method{} is currently evaluated on a separate $128$-source held-out panel with $K=12$, whereas published rows use the official $800$ sources with $K=20$, so starred values are shown as benchmark context rather than a direct numerical ranking. $^{\dagger}$GrIDDD reports diversity without specifying its evaluator implementation, whereas \method{} uses the released Jin evaluator. Diversity details and the broader historical comparison are reported in \cref{app:alignment}.}
%% AUTHORING NOTE, not for the PDF: when the frozen official run lands, replace
%% the starred COMPOSE row in place with (n=800, K=20) and drop the $^{*}$ clause,
%% since the protocols then match.  Keep the diversity column and the dagger.
\label{tab:qed}
\end{table}

%% AUTHORING NOTE, not for the PDF: the result slot below currently holds the
%% prospective held-out panel (n=128, K=12).  When the frozen official run lands,
%% replace the numbers with (n=800, K=20), delete the final sentence because the
%% protocols then match, and rewrite the interpretation to what the result
%% actually supports rather than mechanically asserting competitiveness.
%% KEEP the diversity column and the diversity clause.  Do not drop them merely
%% because success has become directly comparable: success plus diversity is the
%% package GrIDDD reports side by side, and is a better result than success alone.
\method{} achieves strong source-conditioned editing performance under the standard similarity-constrained QED criterion. On the held-out evaluation, \method{} solves $70/128$ sources ($54.7\%$) within $12$ returned candidates and maintains substantial within-source molecular diversity ($0.393$ mean pairwise Tanimoto distance, \cref{fig:capabilities}a), indicating that benchmark success is not driven by collapse to a narrow set of edits. Success increases steadily with candidate allowance rather than depending on a single operating point. \Cref{tab:qed} places this result alongside published methods evaluated under the same benchmark criterion, including GrIDDD at $45.1\%$ success. Because the current \method{} evaluation uses a different source cohort and candidate allowance, these values are reported as benchmark context rather than a direct numerical ranking. Internal inference costs are not matched across methods, since the benchmark constrains the number of returned molecules rather than the successor evaluations performed to produce them.

%% TODO [size-fixed ablation, not yet run] --------------------------------------
%% Run: identical frozen controller, R_theta, sources, candidate allowance and
%% randomness as the 128-source prospective panel above; remove every successor
%% with Delta n != 0.  Motivation: GrIDDD reports the same causal ablation and
%% sees QED success fall 45.1% -> 33.8% (Ninniri et al., NeurIPS 2025, Sec. 5.3).
%% When measured, add as paragraph 3:
%%
%% QED is particularly informative for \method{} because successful optimization
%% can require changes in molecular size. We therefore repeat the evaluation
%% after removing every successor that changes heavy-atom count, while holding
%% $R_\theta$, the controller, sources, candidate allowance, and randomness
%% fixed. Full \method{} solves XXX sources compared with XXX under size-fixed
%% support, a paired difference of XXX percentage points (\cref{fig:capabilities}b).
%% A positive gap would show that trans-dimensional support contributes directly
%% to benchmark performance rather than merely enlarging the formal state space.
%% -----------------------------------------------------------------------------

%% AUTHORING NOTE, not for the PDF: this subsection title asserts a result that
%% has NOT been measured.  Every value below is \XXX.  If the immediate-vs-
%% future-aware contrast comes back neutral or negative, the title must change
%% before submission; do not let a declarative title outlive its measurement.
\subsection{Future-aware control improves Pareto exploration under competing objectives}
\label{sec:exp-pareto}

\Cref{sec:exp-lookahead} showed that remaining-budget information can improve molecular decisions when the desired endpoint is known. We now ask whether the same principle improves multi-objective design when no target molecule is specified. We optimize JNK3 and GSK3$\beta$ across five preference directions following the InversionGNN setting \citep{niu2025inversiongnn}. The immediate and future-aware \method{} controllers share the same frozen $R_\theta$, objective-blind initialization bank, property model, executable edit sets, preferences, and evaluation budget. They differ only in whether a candidate is valued by its properties at the next molecule, $g_\lambda(y)$, or by the preference-weighted futures reachable from it, $h_\phi(y,\lambda,b-1)$. Average Pareto score (APS) is the primary endpoint, while hypervolume, novelty, and diversity are reported in \cref{tab:pareto} as complementary measures of front quality.

\begin{table}[!ht]
\centering
{\small
\begin{tabular}{lrrrr}
\toprule
Method & APS ($\uparrow$) & HV ($\uparrow$) & Novelty ($\uparrow$) & Top-$K$ diversity ($\uparrow$) \\
\midrule
InversionGNN, matched rerun & \XXX{} & \XXX{} & \XXX{} & \XXX{} \\
\method{}, immediate preference & \XXX{} & \XXX{} & \XXX{} & \XXX{} \\
\textbf{\method{}, future-aware} & \textbf{\XXX{}} & \textbf{\XXX{}} & \textbf{\XXX{}} & \textbf{\XXX{}} \\
\bottomrule
\end{tabular}}
\caption{\small\textbf{Two-objective JNK3/GSK3$\beta$ optimization from a common frozen initialization bank.} All rows use the same evaluation budget and preference set. Published InversionGNN results obtained under a different initialization protocol are reported as context in \cref{app:bench-moo} and are not pooled with the matched rerun.}
\label{tab:pareto}
\end{table}

Future-aware control achieves an APS of \XXX{}, compared with \XXX{} under immediate control, an improvement of \XXX{} (\cref{tab:pareto}, \cref{fig:capabilities}b). Hypervolume changes from \XXX{} to \XXX{}, while novelty and diversity remain \XXX{}, indicating {\color{red}[if supported: that the improvement is not obtained by collapsing onto a narrow region of molecular or objective space]}. Because both \method{} arms see the same candidate molecules and identical property predictions, this comparison isolates the effect of propagating preference through the molecular futures induced by $R_\theta$. The gain therefore extends the mechanism of \cref{sec:exp-lookahead} from known-target recovery to target-free multi-objective design.

Under the same initialization, preference set, and evaluation budget, the matched InversionGNN rerun obtains an APS of \XXX{} (\cref{tab:pareto}), providing an external measure of multi-objective task competence. This cross-method comparison is not used to attribute the difference to future-aware control, since the methods do not share a learned property predictor or internal search procedure. These results show that the same frozen molecular process can be redirected toward distinct property trade-offs. We next ask whether this flexibility persists after molecular optimization is already underway and the objective itself changes.

\subsection{\method{}'s realized molecular states preserve progress when objectives change}
\label{sec:exp-retarget}

\Cref{sec:exp-pareto} showed that the same frozen molecular process can be steered toward different property preferences when those preferences are specified in advance. Design priorities, however, can change after optimization is already underway, for example as new measurements introduce liabilities or additional requirements. \method{} is naturally suited to this setting because its reference process $R_\theta$ is goal-independent and every realized state is itself a complete molecule. A newly specified objective can therefore be applied directly to the molecule already reached. We ask whether this capability is practically useful. Does molecular progress accumulated before an objective change provide a better starting point than beginning again from the original molecule?

We test this on $40$ source-disjoint held-out molecules under potency-first and developability-first histories. Each source is edited for three steps before the eventual joint objective $B$ is specified, so the committed prefix cannot anticipate the new goal. After the switch, continuing from the realized state $x_3$ and restarting from the original molecule $x_0$ receive the same $B$, the same three-edit remaining budget, and the same frozen molecular process. This matched comparison isolates whether the molecular state reached before the switch retains useful progress for the newly specified objective.

\begin{table}[!ht]
\centering
{\small
\begin{tabular}{lrr}
\toprule
Comparison & Potency-first & Developability-first \\
\midrule
\textbf{State reuse} (continue $x_3$ vs.\ restart $x_0$) & \textbf{$+0.302$} $[0.186,0.413]$ & \textbf{$+0.176$} $[0.069,0.283]$ \\
Retargeting (new vs.\ old objective) & $+0.748$ $[0.608,0.891]$ & $+1.462$ $[1.240,1.691]$ \\
\bottomrule
\end{tabular}}
\caption{\small\textbf{Reusing realized molecular states after an objective switch.} Entries are paired changes in the post-switch objective, and positive values favor the first strategy. State reuse is the primary comparison, and intervals are source-level bootstrap intervals. The complete record is given in \cref{app:exp-retarget}.}
\label{tab:retarget}
\end{table}

Continuing from $x_3$ improves the post-switch objective over restarting from $x_0$ by $0.302$ $[0.186,0.413]$ after potency-first optimization and $0.176$ $[0.069,0.283]$ after developability-first optimization (\cref{tab:retarget}, \cref{fig:capabilities}c). The benefit is not universal across sources, indicating that state reuse is an empirical advantage rather than a consequence of the controller construction. Retargeting also substantially outperforms continuing under the previous objective, confirming that the subsequent trajectory responds to the newly specified goal. Thus, when design priorities change, \method{} can redirect molecular progress already realized under an earlier objective rather than discarding that progress and beginning again from the original molecule.

%% AUTHORING NOTE, not for the PDF.  This subsection is written in the SUCCESS
%% branch of a confirmation that has NOT been opened.  Numbers render from the
%% \Pw* macros in main.tex, which currently hold DEVELOPMENT (n=24) values from
%% the verified/future-aware arm.  Those values would satisfy both gates, but
%% development did NOT pass the confirmatory test: it is a different, already-
%% inspected sample, and the cLogP corridor was CHOSEN after a feasibility
%% screen.  Never write, or imply, that development passed confirmation.
%% Frozen confirmation: n=48, corridor C=[2.3689,4.4522], intersection-union on
%%    (i)  lower CI bound on hidden-path prevalence  >  1/3
%%    (ii) lower CI bound on utility difference      > -0.25 IQR  (-0.20 sens.)
%% ON THE RESULT:
%%   both pass       -> keep this title and prose; swap the macros; set
%%                      \PwEvalPhrase to "On the 48-source held-out evaluation".
%%   utility fails   -> retitle "Endpoint feasibility can hide pathwise
%%                      molecular constraint violations"; the phenomenon becomes
%%                      the result.  Delete paragraph 4 from "Thus the pathwise
%%                      guarantee" onward, and the "prevents ... while retaining
%%                      comparable endpoint utility" claim with it.
%%   prevalence fails-> the section loses its motivation; compress or demote it
%%                      rather than rescuing it with secondary analyses.
%% Do not change the corridor, either threshold, the controller, or the analysis
%% after seeing the confirmation.  Do not rewrite the question after the answer.
\subsection{\method{} prevents constraint violations hidden by endpoint-only molecular optimization}
\label{sec:exp-pathwise}

\Cref{sec:exp-retarget} showed that realized molecular states can remain useful when design objectives change. Their explicit representation provides a second advantage, in that constraints can govern which molecular states the optimization process is allowed to visit, rather than only which molecule it ultimately returns. This distinction matters whenever intermediate designs must remain within a predictor's applicability domain, preserve a required molecular feature, or satisfy a prescribed physicochemical range. Endpoint-only optimization cannot certify such a requirement from the final molecule alone, whereas \method{} can impose it directly by restricting the executable successor support at every realized state.

We test whether this distinction is consequential using a fixed cLogP corridor $C$ while optimizing an independent DRD2 objective over six edits. Endpoint-only control requires only the returned molecule to lie within $C$, whereas \method{}'s pathwise controller removes successors outside $C$ before a transition is sampled. The two arms otherwise share the same $R_\theta$, sources, objective, horizon, and evaluation budget. Because zero pathwise violations follow from support restriction by construction, the empirical question is not whether the mask works, but whether endpoint-only optimization actually hides meaningful violations that make the stronger guarantee worth having.

Endpoint-only feasibility conceals substantial pathwise violations. \PwEvalPhrase{}, $\PwPrevN{}/\PwPrevD{}$ ($\PwPrevPct{}\%$) of endpoint-feasible trajectories leave $C$ at least once before returning. These excursions extend a median $\PwDepth{}$ cLogP units beyond the admissible region and persist for \PwDuration{} of \PwHorizon{} committed edits (\cref{fig:capabilities}d). Endpoint evaluation can therefore certify a final molecule while overlooking violations across much of the optimization trajectory.

\method{} prevents these hidden violations by restricting the molecular states reachable during optimization. Under future-aware pathwise control, normalized terminal DRD2 utility changes by only $\PwDelta{}$ $\PwCI{}$ relative to endpoint-only control, remaining within the preregistered tolerance. Thus the pathwise guarantee is not merely formal, since \method{} can enforce a requirement throughout molecular optimization while retaining comparable endpoint utility. This capability follows from controlling an explicit process over complete molecular states, in which constraints can act on the trajectory itself rather than being checked only after generation is complete. The complete developmental record, the support-viability audit and the confirmation preregistration are given in \cref{app:exp-pathwise}.

